\begin{mistake}{2026-05-07}{06}

\begin{problemcontent}
设函数 \(f(x)\) 在闭区间 \([a,b]\) 上连续，在开区间 \((a,b)\) 内可导，且 \(f'(x)>0\)。若极限
\[
\lim_{x\to a^+}\frac{f(2x-a)}{x-a}
\]
存在。已知第(I)问可证：在 \((a,b)\) 内 \(f(x)>0\)。证明第(II)问：在 \((a,b)\) 内存在点 \(\xi\)，使
\[
\frac{b^2-a^2}{\int_a^b f(x)\,dx}=\frac{2\xi}{f(\xi)}.
\]
\end{problemcontent}

\begin{solutioncontent}
由第(I)问结论，
\[
f(x)>0,\qquad x\in(a,b).
\]
因此
\[
\int_a^b f(x)\,dx>0,
\]
且对任意 \(\xi\in(a,b)\)，有 \(f(\xi)>0\)，后续分式有意义。

本题使用柯西中值定理。柯西中值定理的形式为：若 \(F,G\) 在 \([a,b]\) 上连续，在 \((a,b)\) 内可导，且 \(G'(x)\neq 0\)，则存在 \(\xi\in(a,b)\)，使
\[
\frac{F(b)-F(a)}{G(b)-G(a)}=\frac{F'(\xi)}{G'(\xi)}.
\]

观察目标式
\[
\frac{b^2-a^2}{\int_a^b f(x)\,dx}=\frac{2\xi}{f(\xi)},
\]
其中 \(b^2-a^2\) 应看成函数 \(x^2\) 的端点差，\(2\xi\) 应看成 \(x^2\) 在 \(\xi\) 处的导数；\(\int_a^b f(x)\,dx\) 应看成变上限积分函数的端点差，\(f(\xi)\) 应看成该变上限积分函数在 \(\xi\) 处的导数。

令
\[
F(x)=x^2,
\qquad
G(x)=\int_a^x f(t)\,dt.
\]

因为 \(F(x)=x^2\) 在 \([a,b]\) 上连续，在 \((a,b)\) 内可导，且
\[
F'(x)=2x.
\]

又因为 \(f\) 在 \([a,b]\) 上连续，所以
\[
G(x)=\int_a^x f(t)\,dt
\]
在 \([a,b]\) 上连续，在 \((a,b)\) 内可导，并且由变上限积分求导公式，
\[
G'(x)=f(x).
\]

由第(I)问，\(f(x)>0\)，故
\[
G'(x)=f(x)>0,
\qquad x\in(a,b),
\]
所以
\[
G'(x)\neq 0.
\]

因此 \(F,G\) 满足柯西中值定理条件，存在 \(\xi\in(a,b)\)，使
\[
\frac{F(b)-F(a)}{G(b)-G(a)}=\frac{F'(\xi)}{G'(\xi)}.
\]

分别计算各部分：
\[
F(b)-F(a)=b^2-a^2,
\]
\[
F'(\xi)=2\xi,
\]
\[
G(b)-G(a)
=
\int_a^b f(t)\,dt-\int_a^a f(t)\,dt
=
\int_a^b f(t)\,dt,
\]
\[
G'(\xi)=f(\xi).
\]

代入柯西中值定理所得等式，得
\[
\frac{b^2-a^2}{\int_a^b f(t)\,dt}
=
\frac{2\xi}{f(\xi)}.
\]
积分变量名称不影响定积分值，故也可写为
\[
\boxed{\frac{b^2-a^2}{\int_a^b f(x)\,dx}=\frac{2\xi}{f(\xi)}}.
\]

补充说明：若改用积分中值定理，只能得到存在 \(\eta\in(a,b)\)，使
\[
\int_a^b f(x)\,dx=(b-a)f(\eta).
\]
这里的 \(\eta\) 与第(II)问中的 \(\xi\) 没有理由相同，因此不能把 \(\eta\) 直接写成 \(\xi\)，也不能由此推出 \(\xi=(a+b)/2\)。柯西中值定理的优势是一次性给出同一个点 \(\xi\)，使 \(2\xi\) 和 \(f(\xi)\) 同时出现。
\end{solutioncontent}

\mistakeknowledge{
\begin{itemize}
  \item \textbf{核心公式/定理}：柯西中值定理：端点差之比 \(\frac{F(b)-F(a)}{G(b)-G(a)}\) 可转化为同一点处的导数之比 \(\frac{F'(\xi)}{G'(\xi)}\)。本题取 \(F(x)=x^2\)，\(G(x)=\int_a^x f(t)\,dt\)。
  \item \textbf{易错点}：积分中值定理给出的点 \(\eta\) 与第(II)问要证的点 \(\xi\) 不一定相同，不能把 \(\int_a^b f(x)\,dx=(b-a)f(\eta)\) 中的 \(\eta\) 偷换成 \(\xi\)。
  \item \textbf{关联知识}：看到 \(b^2-a^2\) 与 \(2\xi\)，联想 \(F(x)=x^2\)；看到 \(\int_a^b f(x)\,dx\) 与 \(f(\xi)\)，联想 \(G(x)=\int_a^x f(t)\,dt\)。
  \item \textbf{审题条件}：第(I)问的 \(f(x)>0\) 用于保证 \(G'(x)=f(x)\neq 0\)，从而满足柯西中值定理条件，并保证分母 \(f(\xi)\) 与 \(\int_a^b f(x)\,dx\) 不为零。
\end{itemize}}

\end{mistake}
